\documentclass[12pt]{article}
\usepackage[margin=1in]{geometry}
\usepackage{amsmath,amssymb,amsthm,amsfonts}
\usepackage{hyperref}
\hypersetup{colorlinks=true, linkcolor=blue, urlcolor=blue, citecolor=blue}

\theoremstyle{definition}
\newtheorem*{definition}{Definition}
\theoremstyle{plain}
\newtheorem*{theorem}{Theorem}
\newtheorem*{lemma}{Lemma}

\title{\textbf{Appendix: Detailed Proof of $\Pi^0_2$-Hardness for Lossless Channels}}
\author{Dmytro Surdu}
\date{}

\begin{document}
\maketitle

\section*{Necessity of Information Loss via Reduction from \textsc{TOTAL}}

\begin{definition}[The \textsc{TOTAL} Problem]
Let $\{M_e\}_{e\in\mathbb N}$ be a standard enumeration of Turing machines.  The set
\[
\mathrm{TOTAL} \;=\;\bigl\{\,e\in\mathbb N:\;\forall x\in\mathbb N\;\exists t\;M_e(x)\text{ halts within }t\text{ steps}\bigr\}
\]
is $\Pi^0_2$–complete.
\end{definition}

We now show that for any lossless (invertible) discrete measurement channel, the corresponding infinite–tail certification problem is at least as hard as \textsc{TOTAL}.  In particular, no finite deterministic witness can collapse this to an $\mathbf{NP}$ decision problem.

\begin{definition}[Discrete Lossless Measurement Model]
We work in the following extremely simple, purely \emph{discrete} model:
\begin{itemize}
  \item \textbf{State space:} $\mathcal Q=\mathbb N$.
  \item \textbf{Channel:} $g\circ f : \mathcal Q\to\mathcal Q$ is the identity map, which is a bijection (hence lossless).
  \item \textbf{Stream generator $h$:} For each initial state $q=e\in\mathbb N$, $h(e)$ outputs an infinite sequence over the alphabet $\{1,\#, \bot\}$ by dovetailing simulations of $M_e(1),M_e(2),\dots$ as follows:
    \begin{enumerate}
      \item Interleave one step of $M_e(1)$, then one step of $M_e(2)$, then one step of $M_e(1)$, then one step of $M_e(3)$, and so on, ensuring every $M_e(x)$ makes progress.
      \item Whenever $M_e(x)$ is discovered to halt in $t$ steps, immediately append
        \[
          1^x\,\#\,1^t\,\#
        \]
        to the output stream, then resume the dovetailed simulation.
      \item If $M_e(x)$ never halts, no certificate for that $x$ ever appears.
      \item After every halting event the process continues indefinitely, so the overall output is infinite.
    \end{enumerate}
\end{itemize}
\end{definition}

\begin{definition}[Tail Predicate \(P_e\)]
For each $e\in\mathbb N$, define the infinite–tail predicate $P_e$ on streams $Z\in\{1,\#,\bot\}^\omega$ by
\[
  P_e(Z)\quad:\quad
  \forall x\in\mathbb N\;\exists t\in\mathbb N\text{ s.t.\ }1^x\#1^t\#\text{ appears as a contiguous substring of }Z.
\]
\end{definition}

\begin{theorem}
Under the discrete lossless model above, deciding whether a given stream satisfies $P_e$ is $\Pi^0_2$–hard.  Hence, without information loss (invertibility), the certification problem lies outside $\mathbf{NP}$ unless the Polynomial Hierarchy collapses.
\end{theorem}

\begin{proof}
We give a many–one reduction from \textsc{TOTAL} to the certification problem for $P_e$.

\medskip\noindent\textbf{(i) Correctness of the reduction.}
Let $e\in\mathbb N$ be arbitrary.  Consider the stream
\[
  Z \;=\; h(e)\;=\;h_e(e).
\]
By construction:

\begin{itemize}
  \item[$\Rightarrow$] If $e\in\mathrm{TOTAL}$, then for every $x\in\mathbb N$, the machine $M_e(x)$ halts in some finite number of steps $t_x$.  During the dovetailed simulation of all $M_e(x)$, each halting event is eventually discovered, and so for each $x$ the certificate block $1^x\#1^{t_x}\#$ is appended to $Z$.  Consequently $\forall x\,\exists t\,\bigl(1^x\#1^t\#\text{ appears in }Z\bigr)$, i.e.\ $P_e(Z)$ holds.

  \item[$\Leftarrow$] Conversely, if $P_e(h(e))$ holds, then for every $x$ there is some $t$ such that $1^x\#1^t\#$ appears in the output.  By the definition of $h$, the only way such a block can ever be produced is upon discovering that $M_e(x)$ indeed halts in exactly $t$ steps.  Hence $M_e(x)$ halts for all $x$, and so $e\in\mathrm{TOTAL}$.
\end{itemize}

Thus
\[
e\in\mathrm{TOTAL}
\;\Longleftrightarrow\;
P_e\bigl(h(e)\bigr).
\]

\medskip\noindent\textbf{(ii) Complexity implications.}
Because $g\circ f$ is the identity on $\mathbb N$, a finite observation of the stream’s “head” gives \emph{exactly} the initial state $e$.  There is no loss of information at the channel stage.  Any deterministic verifier that, given a finite witness and a finite prefix of the stream, tries to decide $P_e$ must effectively solve
\[
  \forall x\;\exists t\;M_e(x)\text{ halts in }t\text{ steps},
\]
which is a canonical $\Pi^0_2$–complete problem.  No finite, polynomial–size witness can encode “all future halts” into a single certificate checkable in polynomial time.  Hence the certification language cannot lie in $\mathbf{NP}$ unless the Polynomial Hierarchy collapses.

\end{proof}

\end{document}